\mysection{実行可能ファイルのパッチ}

\subsection{テキスト文字列}

C文字列は（暗号化されていない限り）、どのヘキサエディターでもパッチするのが最も簡単なものです。
この技術は、機械語コードや実行可能ファイル形式を知らない人でも利用できます。
新しい文字列は古い文字列より大きくてはいけません。なぜなら、そこで別の値やコードを上書きするリスクがあるからです。
\myindex{MS-DOS}

この方法を使用して、80年代と90年代の旧ソ連諸国では、少なくともMS-DOS時代に多くのソフトウェアが\IT{ローカライズ}されました。
\IT{ローカライズ}されたソフトウェアに奇妙な略語が存在していた理由：より長い文字列のためのスペースがなかったからです。

\myindex{Borland Delphi}

Delphi文字列については、必要に応じて文字列のサイズも修正する必要があります。

\subsection{x86コード}
\label{x86_patching}

頻繁なパッチタスク：

\myindex{x86!\Instructions!NOP}
\begin{itemize}

\item 
最も頻繁な作業の一つは、何らかの命令を無効にすることです。
これは多くの場合、バイト\TT{0x90}（\ac{NOP}）で埋めることによって行われます。

\item \TT{74 xx}（\JZ）のようなオペコードを持つ条件ジャンプは、2つの\ac{NOP}で埋めることができます。

また、2番目のバイト（\IT{ジャンプオフセット}）に0を書き込むことで条件ジャンプを無効にすることも可能です。

\myindex{x86!\Instructions!JMP}
\item 
もう一つの頻繁な作業は、条件ジャンプを常にトリガーするようにすることです：
これは、\JMPを表す\TT{0xEB}をオペコードの代わりに書き込むことで実行できます。

\myindex{x86!\Instructions!RET}
\myindex{stdcall}
\item 関数の実行は、その開始時に\RETN（0xC3）を書き込むことで無効にできます。
これは\TT{stdcall}（\myref{sec:stdcall}）を除くすべての関数に当てはまります。
\TT{stdcall}関数をパッチする際は、引数の数を決定し（例えば、この関数内で\RETNを見つけることで）、
16ビット引数を持つ\RETNを使用する必要があります（0xC2）。

\myindex{x86!\Instructions!MOV}
\myindex{x86!\Instructions!XOR}
\myindex{x86!\Instructions!INC}
\item 時々、無効にされた関数は0または1を返す必要があります。
これは\TT{MOV EAX, 0}または\TT{MOV EAX, 1}で行うことができますが、少し冗長です。\\
より良い方法は\TT{XOR EAX, EAX}（2バイト\TT{0x31 0xC0}）または\TT{XOR EAX, EAX / INC EAX}（3バイト\TT{0x31 0xC0 0x40}）です。

\end{itemize}

ソフトウェアは変更に対して保護されている可能性があります。

この保護は多くの場合、実行可能コードを読み取ってチェックサムを計算することによって行われます。
したがって、保護がトリガーされる前にコードを読み取る必要があります。

これは、メモリ読み取りにブレークポイントを設定することで判断できます。

\myindex{tracer}
\tracerには、このためのBPMオプションがあります。

PE実行可能ファイルのリロケーション（\myref{subsec:relocs}）は、パッチ中に触れてはいけません。
なぜなら、Windowsローダーが新しいコードを上書きする可能性があるからです。
\myindex{Hiew}
（例えば、Hiewではグレー表示されています：
\figref{fig:scanf_ex3_hiew_1}）。

最後の手段として、リロケーションを回避するジャンプを書くか、リロケーションテーブルを編集する必要があります。